% Seminar paper, Seminar Urban Computing, Human-IST, University of Fribourg (AS-2026).
% ACM acmart, single-column manuscript, author-year citations (Human-IST seminar guidelines).
% Build: see paper/README.md. Every number comes from numbers.tex (invariant I6):
% never type a measured value here; use the macros, or put tables in paper/tables/.
\documentclass[manuscript,screen,nonacm]{acmart}
\citestyle{acmauthoryear}

\usepackage{xcolor}
\newcommand{\placeholder}[1]{\textcolor{gray}{[#1]}}
% A number quoted from cited literature (not from results.json), e.g.
% "the operator reports \citednum{94\%} punctuality", next to its citation.
% The I6 checker skips it. Never use it for our own results.
\newcommand{\citednum}[1]{#1}

\input{numbers}  % generated: make paper-numbers (scripts/export_paper_numbers.py)

\begin{document}

\title{Can I Trust My Connection? Fuzzy Linguistic Summaries as the Explainable Core of a
Swiss Public-Transport Assistant}

\author{Nikhil Shetty}
\affiliation{%
  \institution{University of Fribourg, Human-IST Institute}
  \city{Fribourg}
  \country{Switzerland}}

% UC-072 writes the final abstract (about 150 words) once the results exist.
\begin{abstract}
Commuters ask simple questions about their own trips: is my connection reliable, and how
late is my train usually? Official punctuality figures answer with a crisp rule and
network-wide averages, while chatbots answer with vague words that nobody checks. We
present a citizen-facing assistant for Swiss public transport whose explanations are fuzzy
linguistic summaries computed from open operational data. The decision word shown to the
passenger and its explanation come from the same truth computation, so the explanation is
faithful by construction. Using the quality measures of Hudec and colleagues, we test which
summaries still hold for a passenger's own line, station and time of day, and how often
network-wide statements hide local reality. We compare the official crisp view with a view
based on the passenger's own transfer margin. Finally, we audit a local large language
model that phrases such answers, scoring its quantified claims with a crisp and a fuzzy
checker. \placeholder{Key results with numbers from numbers.tex, UC-072.}
\end{abstract}

\keywords{computing with words, fuzzy linguistic summaries, explainability, public
transport, open government data, large language models}

\maketitle
\ResultsBanner

% UC-071: Introduction (problem, research questions, contribution).
\section{Introduction}
\label{sec:intro}
\placeholder{UC-071: problem statement, the main research question and SRQ1--SRQ3,
contributions, paper outline.}

% UC-071: Background and related work. Cite only entries verified in docs/literature.md.
% Sources marked "verified (metadata)" are cited only for what docs/literature.md records
% next to them (e.g. "as attributed by ...").
% TODO(verify): add verified entries for Zadeh's computing with words and Mendel's
% Perceptual Computer to docs/literature.md, then cite them in Section 2.1.
% Pincay 2022 (issue #53 open): claim only what pincay2021 §5 states.
\section{Background and Related Work}
\label{sec:related}

The assistant draws on four strands of work: computing with words and its use in citizen
interfaces, fuzzy linguistic summaries and their quality measures, explanation through
linguistic summaries, and the verification of quantified statements produced by language
models. A fifth, smaller strand concerns how passengers perceive punctuality.

\subsection{Computing with Words in Citizen-Facing Systems}
\label{sec:rw-cww}
Computing with words (CWW) treats words as the objects of computation: their meaning is a
fuzzy set, and a system reasons with these meanings instead of with crisp categories. A
common way to structure such a system is the architecture of the Perceptual Computer, in
which an \emph{encoder} turns inputs into fuzzy sets, a \emph{CWW engine} computes with
them, and a \emph{decoder} turns the result back into words a person can act on. Our
assistant follows this structure (Section~\ref{sec:concept}).

At Human-IST, CWW has been studied as a way to make interfaces respond to people's
perceptions rather than to rigid categories: \citet{colombo2024} proposes \emph{phenotropic
interaction}, interfaces improved with computing with words and perceptions. The limits of
crisp categories in public data are the subject of \citet{spycher2026}, who examine
taxonomies for open government data. Our work shares this perspective: the official
punctuality rule is a crisp category over open government data, and a passenger's notion of
``late enough to miss my train'' is a perception that depends on their situation.

\subsection{Fuzzy Linguistic Summaries and Their Quality}
\label{sec:rw-ls}
A linguistic summary is a quantified sentence about a data set, such as ``most trains are on
time''. The idea goes back to \citet{yager1982}. \citet{kacprzyk2009} describe such
summaries through \emph{protoforms}, templates such as ``$Q$ of the $x$ are $S$'' and ``$Q$ of
the $R$ $x$ are $S$'', in which $Q$ is a fuzzy quantifier (``most''), $S$ a summarizer
(``on time'') and $R$ a restriction (``in the morning peak''). The truth, or validity, of a
summary is the membership of the summarizer's proportion in the quantifier; we use it in
the form given by \citet[Eq.~(2)]{pincay2021}, who attribute it to \citet{kacprzyk2009}, and
by \citet[Eqs.~(5)--(6)]{hudec2022}, who attribute it to \citet{yager1982}.

Validity alone does not say how much data a sentence rests on. \citet{hudec2018} use short
quantified sentences to disseminate official statistics and combine validity with
\emph{coverage}: a restricted summary may be fully valid and yet describe only a handful of
records. Their quality measure $Q_c$ is a non-idempotent t-norm of validity and coverage
(for example the product, which we use) when the coverage reaches a threshold, and zero
otherwise \citep[Eq.~(9)]{hudec2018}; they attribute
the merge of validity and coverage to \citet{hudec2017}. The degree of sufficient coverage is
a smooth function of the share of records that support both the restriction and the
summarizer \citep[Eqs.~(A.3)--(A.4)]{hudec2018}, a function \citet{hudec2022} attribute to
\citet{wu2010}. \citet{hudec2022} further propose a \emph{regional} quality measure for
summaries shown on maps: the same sentence is evaluated per region, and its truth is combined
with the region's relative size through an ordinal sum of a conjunctive and a disjunctive
operator, so that small regions are neither ignored nor allowed to dominate. These are the
measures we use for SRQ1.

\subsection{Linguistic Summaries of Traffic Data}
\label{sec:rw-traffic}
\citet{pincay2021} mine linguistic summaries from road traffic data with the quantifiers
\emph{few}, \emph{about half} and \emph{most}, a validity threshold for reporting a summary,
and the Simpful library for fuzzy logic \citep{spolaor2020}. In their conclusion they name
two extensions as future work: combining validity with existing quality measures in the
sense of Hudec et al., and checking validity by region, to see whether the data distribution
skews the summaries \citep[§5]{pincay2021}. The same group has continued with geospatial
traffic data \citep{pincay2022} and with Swiss traffic safety analysis combining machine
learning and a local language model \citep{pincay2026}. Our work transfers the approach to
public transport and implements both extensions named in \citet{pincay2021}; we use their
configuration as a methodological baseline, not as a reproduction of their numbers, since
their data is not public.

\subsection{Linguistic Summaries as Explanations}
\label{sec:rw-xai}
Linguistic summaries have been proposed as a form of explainable artificial intelligence.
\citet{plenary2022} explain the behaviour of black-box models in natural language through
fuzzy linguistic summaries, and \citet{wrede2023} work towards making linguistic summaries
themselves explainable. In PLENARY the summaries describe a separate black-box model; in our
assistant the summaries are the model: the decision word is decoded from
the same fuzzy proportion that the main explanation sentence states, so the explanation
cannot diverge from how the decision was made.

\subsection{Language Models and Quantified Statements}
\label{sec:rw-llm}
\citet{eida2026} study how language models generate quantified statements from tables and
verify them by executing checks against the source table. Their truth conditions are crisp
and set-based; ``most'' is true if more than half of the relevant rows satisfy the property
\citep{eida2026}. Their tables are synthetic and small, and their verification
functions are themselves written by a language model. We adopt their set-based truth
conditions as our crisp checker, implemented as deterministic code, and add a fuzzy checker
that assigns graded truth to vague quantifiers and summarizers, on real data.
\citet{gaitan2025} combine language models with fuzzy linguistic summaries of raw sensor
data in diabetes care, with a methodology for fine-tuning and evaluating the models; the
truth of the quantifiers that the language model writes is not recomputed from the data. The
language model in our audit is Apertus \citep{apertus2025}, an open model from the Swiss AI
Initiative that we run locally.

\subsection{Passengers' Perception of Punctuality}
\label{sec:rw-perception}
Official punctuality measures and passengers' experience differ. A study of passengers in
Great Britain found that passengers judge punctuality more strictly than the industry
measure and expect trains to arrive within about one minute of the timetable
\citep{transportfocus2015}. We use this finding as motivation only. Swiss perception data
does not exist for our purpose, and we do not collect new data from people. Instead, our
perception-aware view makes ``late'' depend on a quantity the passenger knows precisely:
their own time to change trains.

\subsection{Positioning}
\label{sec:rw-position}
Taken together, linguistic summaries with quality measures exist for official statistics
and maps \citep{hudec2018,hudec2022}, linguistic summaries of traffic exist without them
\citep{pincay2021}, linguistic summaries explain other models \citep{plenary2022} and are
themselves made explainable \citep{wrede2023},
and quantified statements of language models are verified crisply \citep{eida2026}. This
paper connects these strands in one citizen-facing artefact: summaries with Hudec's quality
measures on public transport data are the decision core of an assistant, the decision depends
on the passenger's own margin, and the same truth computation serves as a fuzzy checker for
language-model answers, compared against a crisp one.

% UC-071: Concept part (CWW assistant design).
\section{Concept: A Computing-with-Words Assistant}
\label{sec:concept}
\placeholder{UC-071: vocabulary, truth of summaries, quality measures, decision rule and
faithful explanation (docs/methods.md sections 1--6).}

\subsection{Vocabulary and Truth}
\placeholder{UC-071.}

\subsection{Decision and Explanation}
\placeholder{UC-071.}


\subsection{Methodological Baseline}
\label{sec:methodological-baseline}
We apply the unrestricted linguistic-summary truth computation of
\citet[Eq.~2]{pincay2021} as a methodological baseline for Swiss train arrival delays.
The input is the study's measured Ist-Daten arrivals, without a passenger-specific
restriction, and the summarizers are the study's four arrival-delay words: on time,
slightly late, late and very late. For each word, its fuzzy proportion is evaluated
with the quantifiers few, about half and most. We compute these scores for the pooled
network and separately for each study half.

The baseline uses the dedicated \texttt{pincay2021} vocabulary configuration, whose
curves follow \citet[Fig.~2, p.~172]{pincay2021}. The main study retains its own
quantifier breakpoints (D-028). A summary is valid when its truth is at least the
configured mining threshold $\beta$; equality is included. The default follows
\citet[Section~\citednum{3.2}, p.~172]{pincay2021}, while configuration overrides also apply to
this baseline. This is a methodological comparison on Swiss arrival data, not a
reproduction or numerical comparison of the paper's traffic-message results.

% UC-071: Prototype (practical part); installation and manual are in the appendix (UC-061).
% TODO(verify): add a verified design-science reference to docs/literature.md before citing
% one here; the framing below is stated without citation.
\section{Prototype}
\label{sec:prototype}

The prototype is the artefact of this design-oriented study: it makes the concept of
Section~\ref{sec:concept} usable by a commuter and serves as the object of the evaluation.
Its design follows three requirements derived from the problem. First, the answer must be
about the passenger's own trip, not about the network. Second, every statement must be
checkable against the data, so the interface shows how true each sentence is and how much
data it rests on. Third, the assistant must not invent anything, so no language model is
called at run time.

\subsection{Architecture}
\label{sec:architecture}
The prototype is a web application built with Streamlit on top of the same library that
produces the evaluation results. The data pipeline ingests the daily Ist-Daten files,
validates them, builds a table of measured train arrivals and reconstructs the connections
at the three study hubs (Zürich HB, Bern and Olten) from pairs of arriving and departing
trains. Hub stops are matched by their official identifiers rather than by display names, so
the crisp and the fuzzy view always use the same arrivals. The assistant library contains
the vocabulary, the truth engine, the quality measures, the summary miner, the decision
engine and the crisp official view. The user interface only renders what the library
returns: a separate, tested rendering layer derives every text and number on the screen from
the answer object, so the interface cannot show a number the engine did not compute.

\subsection{User Interface}
\label{sec:ui}
The commuter enters the question in a form: the hub, the arriving and the departing line, the
time of day, the days and their time to change, chosen with a slider. Station names and time
bands come directly from the study configuration, so that the questions match the data
exactly. The answer is organised in three views (Appendix~\ref{app:manual} and
Figures~\ref{fig:app-form}--\ref{fig:app-chatbot}).

\subsubsection*{Answer} A coloured badge shows the decision word. Below it, under \emph{Why},
each explanation sentence is shown with a bar for its truth, coloured by how true the
sentence is (green, amber or red), the number of observations it is based on and its
coverage. The colour of the bar describes the sentence, not the risk: a fully true sentence
that explains a risky answer has a green bar next to a red badge.

\subsubsection*{Official view} The passenger's verdict and its sentence are shown side by side
with the official punctuality verdict and its crisp sentences. This view is the
citizen-facing form of SRQ2: when the two verdicts differ, the passenger sees that their own
time to change is what makes the difference.

\subsubsection*{Chatbot view} For questions of the evaluation set, the view shows the answer
that the language model wrote, stored from the offline runs, with a coloured mark per claim
from the fuzzy checker and a grey mark for claims that cannot be checked. This view is the
citizen-facing form of SRQ3.

Every view ends with the data source and the month of the measurements. The assistant is not
branded as a product of any operator, because the data covers all operators.

\subsection{Design Choices}
\label{sec:design}
Several choices follow from the requirement that every statement be checkable.
\begin{itemize}
  \item \emph{One sentence about the passenger, always.} Every answer contains a sentence
  about the passenger's own connections, either with the best-fitting quantifier or, in the
  vocabulary gap, with the measured share. A risky verdict without any statement about the
  passenger's connections would ask the passenger to trust the decision blindly.
  \item \emph{Truth next to the sentence, not hidden.} Showing the truth and the number of
  observations lets a passenger see that ``few of your connections are missed'' rests on many
  trains or on a handful.
  \item \emph{Both views on one screen.} The official verdict is not hidden but shown next to
  the passenger's verdict, with its own sentences, so the passenger can see why the two
  differ.
  \item \emph{No generation at run time.} The chatbot view shows stored answers from the
  evaluation only. A citizen-facing tool that generates new sentences would need the
  checker in the loop for every sentence, which is outside the scope of this prototype.
\end{itemize}

\subsection{Example}
\label{sec:example}
On the synthetic test data shown in the figures, a commuter at Zürich HB arriving with the
IC1 in the morning peak and three minutes to change receives the answer \emph{risky}: about
half of the observed connections would have been missed. With ten minutes to change, the
same trains give the answer \emph{safe}. The official view, which ignores the time to
change, rates the same arrivals \emph{risky}, because about half of them are three or more
minutes late. The figures will be replaced by screenshots of the real data once the full
evaluation has run.

\subsection{Verification in the Prototype}
\label{sec:proto-verification}
Each answer the prototype shows can be verified with the same function that the evaluation
uses as a correctness gate (Section~\ref{sec:explanation}). The rendering layer is tested to
show only numbers that appear in the answer, and the application is tested end to end on the
synthetic data. Installation from a clean machine takes five commands
(Appendix~\ref{app:install}).

% UC-072: Empirical study, part 1 (data and method).
\section{Empirical Study}
\label{sec:study}

\subsection{Data}
\label{sec:data}
\placeholder{UC-072: opentransportdata.swiss actual data, study period, hubs, cleaning
(docs/data.md), invariants. Results provenance: git \texttt{\ResultsGitSha}, config
\texttt{\ResultsConfigHash}.}

% UC-072: Empirical study, part 2 (SRQ1-SRQ3). Numbers ONLY via numbers.tex macros.
\subsection{Results}
\label{sec:results}

\begin{table}[h]
  \caption{Pre-registered hypotheses (eval/hypotheses.yaml): estimate, \CILevel{} confidence
  interval, p-value and verdict, all from results.json.}
  \label{tab:hypotheses}
  \begin{tabular}{llllll}
    \toprule
    RQ & Hypothesis & Estimate & \CILevel{} CI & $p$ & Verdict \\
    \midrule
    SRQ1 & H1a & \HoneaPoint & \HoneaCI & \HoneaP & \HoneaVerdict \\
    SRQ1 & H1b & \HonebValue & -- & -- & \HonebVerdict \\
    SRQ1 & H1c & \HonecPoint & \HonecCI & \HonecP & \HonecVerdict \\
    SRQ2 & H2a & \HtwoaPoint & \HtwoaCI & \HtwoaP & \HtwoaVerdict \\
    SRQ2 & H2b & \HtwobPoint & \HtwobCI & \HtwobP & \HtwobVerdict \\
    SRQ3 & H3a & \HthreeaPoint & \HthreeaCI & \HthreeaP & \HthreeaVerdict \\
    SRQ3 & H3b & \HthreebPoint & \HthreebCI & \HthreebP & \HthreebVerdict \\
    SRQ3 & H3c & \HthreecPoint & \HthreecCI & \HthreecP & \HthreecVerdict \\
    \bottomrule
  \end{tabular}
\end{table}

\begin{table}[h]
  \caption{Methodological baseline for Swiss train arrivals. Each cell is the truth of
  an unrestricted quantifier--delay-word summary. The baseline uses the quantifier
  curves of \citet[Fig.~2, p.~172]{pincay2021}, rather than the main study's own
  breakpoints (D-028); the arrival-delay summarizers are those of the study.
  Validity uses the inclusive configured threshold $\beta$.
  Values await the study-data evaluation; dashes denote pending results, not zero truth.}
  \label{tab:baseline-pincay}
  \begin{tabular}{llll}
    \toprule
    Arrival-delay word & Few & About half & Most \\
    \midrule
    On time & -- & -- & -- \\
    Slightly late & -- & -- & -- \\
    Late & -- & -- & -- \\
    Very late & -- & -- & -- \\
    \bottomrule
  \end{tabular}
\end{table}

Table~\ref{tab:baseline-pincay} reserves the pooled baseline layout. Its values must be
populated from the runner's \texttt{baseline\_pincay.csv}; the corresponding
\texttt{baseline\_pincay\_first.csv} and \texttt{baseline\_pincay\_second.csv} provide
per-half results for the empirical analysis. No study-data baseline values are yet
reported here.

\subsubsection{SRQ1: Is the answer about my trip?}
\placeholder{UC-072.}

\subsubsection{SRQ2: At risk for whom?}
\placeholder{UC-072.}

\subsubsection{SRQ3: Can I trust the chatbot?}
\placeholder{UC-072.}

% UC-072: Discussion and limitations.
\section{Discussion}
\label{sec:discussion}
\placeholder{UC-072.}

\section{Limitations}
\label{sec:limitations}
\placeholder{UC-072: reconstructed connections and minimum transfer time (D-009), one
month of data, one language model family, no user study (D-004), exploratory tests without
multiplicity correction.}

% UC-072: Conclusion.
\section{Conclusion}
\label{sec:conclusion}
\placeholder{UC-072.}

% UC-072: Declaration on the use of generative AI (Human-IST seminar guidelines).
\section*{Declaration on the Use of Generative AI}
\placeholder{UC-072: which tools were used for what (code written by coding agents and
reviewed by the author; paper drafting), and the author's responsibility. Model names
with version numbers go in a macro or a table (invariant I6).}


\bibliographystyle{ACM-Reference-Format}
\bibliography{references}

\appendix
% UC-061: prototype installation guide and user manual (Human-IST guidelines).
\section{Prototype: Installation and User Manual}
\label{app:prototype}
\placeholder{UC-061: installation (docs/INSTALL.md) and user manual (docs/USER\_MANUAL.md).}


\end{document}
